\documentclass[11pt]{article}
\usepackage[utf8]{inputenc}
\usepackage[T1]{fontenc}
\usepackage{geometry}
\usepackage{enumitem}
\usepackage{xcolor}
\usepackage{hyperref}
\usepackage{amsmath}

\geometry{margin=1in}

\title{\textbf{Singular Learning Theory for Applied Deep Learning}\\[8pt]
\large Paper Goals (v5)}
\author{Working Document}
\date{\today}

\begin{document}
\maketitle

\noindent\emph{Companion to the applied-statistics SLT paper: same arc
(urgency $\to$ geometric picture $\to$ formalize $\to$ deliver $\to$
conclude and open a horizon), reorganized for a deep-learning audience.
Open items live in the needs document; this document states only what is
settled, and marks contingent phrasing inline with [contingent: Cn]
tags whose fallbacks live in \S6.}

% ============================================================================
\section{Primary Goal}
% ============================================================================

Deep learning models are singular: the minimum of the loss is not a
point but a positive-dimensional set, and the geometry of the loss
around that set --- manufactured by the architecture and the data
generating process --- governs what the practitioner actually cares
about. Singular learning theory is the mathematics of exactly this
situation: mature, increasingly computable, and still largely unknown
to the practitioners whose models it describes. We write this paper to
communicate SLT to the deep learning practitioner: the correct
geometric picture, its consequences for the tools they use daily, and
the toolkit now being built on it.

This paper aims to:

\begin{enumerate}[nosep, label=(\roman*)]
\item \textbf{Establish the stakes:} before the communication begins,
      show the reader why it is addressed to them --- at two levels.
      Level 1: the set sits under the numbers they routinely report
      (uncertainty, attribution, complexity, selection); the tools that
      assume it away fail at their own stated jobs and mask the failure
      behind plausible numbers, structurally wrong where checkable
      [contingent: C1]. Level 2: the geometry around the set reaches
      their practice even where nothing is wrong --- it is why some
      hyperparameters work and others don't, why some algorithms land
      better than others; training silently selects from the set
      without saying so. Not a theoretical curiosity: their
      architectures, their numbers, their conclusions.
\item \textbf{Correct the mental picture:} replace the flat-valley
      intuition with the singular one --- a set the objective cannot
      choose from and training silently does, and the geometry around
      it: how the loss vanishes along it, how the posterior conforms to
      that structure rather than to any Gaussian, how sample size
      decides what is resolved of it, and the invariant ($\lambda$)
      that summarizes it --- so the reader understands \emph{why} the
      failures occur, \emph{when} to expect them, and \emph{what
      question} each of their numbers actually answers.
\item \textbf{Recast existing tools, provide new ones, and mark what
      is established versus open.} Show why the tools the reader
      already trusts --- cross-validation, Adam, weight decay --- work,
      for singular-geometric reasons the field has misnamed; hand over
      the singular-first toolkit; and be honest about how mature each
      piece is, at the reader's scale.
\item \textbf{Flag the extension of singular geometry beyond the
      loss.} For the same reasons the loss is singular, one should
      expect activation and representation spaces, and policy/value
      spaces, to be singular, and interpretability to be governed by
      degeneracy on the singular set. Explicitly conjectural: one
      developed flag (interpretability), two credible sketches
      (activations, policies).
\end{enumerate}

The reader should finish thinking: ``The minimum I found was a set,
and what I actually care about --- generalization, uncertainty, model
selection --- is governed by the geometry around that set, not by my
parameter count. My tools assumed it away, and my training chose from
it without telling me. Now I understand why my numbers misbehave, what
question each one actually answers, and how to measure the geometry
itself --- and I know where the current limits are, so I won't
overreach.''

% ============================================================================
\section{What the Reader Walks Away With}
% ============================================================================

\begin{enumerate}[label=\textbf{\arabic*.}, leftmargin=*, itemsep=6pt]

\item \textbf{The diagnosis: a missing language, not poor judgment.}
The reader sees that the field has met the phenomena --- mode
connectivity, degenerate Hessians, mergeable checkpoints --- and even
mapped the set they trace out, without a language for the geometry that
governs them: how the loss vanishes around the set, and how the
posterior and sample size answer to that. And the reader can name the
cost of that gap: wrong results, misaimed fixes, years of regularizing
instead of redesigning, and folklore where there should be theory.

\item \textbf{One cause behind the failures, the folklore, and much
that works.}
The scattered failures are not separate bugs, and the familiar
frustrations --- the damping constant no one can justify, sharpness
measures that disagree after a reparameterization, checkpoints that
average into a working model --- are not separate mysteries: one
geometry produces them, explains why the trusted tools succeed at their
jobs, and names what training silently chooses.

\item \textbf{The masking.}
The reader understands why they never noticed: the repairs were aimed
at stability and delivered it; every regularity-restoring repair
returns a plausible number. The failure mode of singular geometry in
deep learning is silence, not error messages.

\item \textbf{The geometric picture without the algebra.}
Taught as a working intuition, in the order a course would: continuous
symmetry and overparameterization force the minimizers to be a
positive-dimensional set; the flat thing the practitioner finds is a
\emph{normal crossing} --- intersecting sheets, hyperbolic
(product-vanishing) level sets --- not a wide valley; the crossing is
built of strata, and the loss vanishes to different orders along and
across them, which is where the consequences live. And the reader can
probe geometry they cannot plot: eigenspectra, transverse slices, the
posterior on the strata.

\item \textbf{The framework, and the lens it grounds.}
Watanabe's central move in a practitioner's language: the parameter
space is stratified, not a manifold; the blow-up transforms it to
normal-crossing form; $\lambda$ replaces $p/2$; integral-based
quantities survive the geometry, local-differential ones do not ---
carried out in full rigor on an analytic microscope model, the ReLU gap
stated, not stepped around. And the habit it installs: three kinds of
hardness (geometric, sample-size, computational), and of any number a
tool returns, ask \emph{which question it actually answers} --- a
lookup, not a procedure.

\item \textbf{Tools graded by the question they answer, and the choice
made measurable.}
Two questions, two criteria: prediction (CV/WAIC, vindicated in
singular models where AIC/DIC provably fail) versus identifying
structure (free energy/WBIC, consistent where CV is not). Minimization
on a singular landscape is an incomplete specification --- the loss
cannot prefer among loss-equivalent points; the optimizer and its
hyperparameters silently complete the choice --- so the trusted tools
are recast, not retired (Adam honest and silent on its choice; weight
decay and flatness re-read), and the LLC reads out where training
landed. The instruments --- singular VI, susceptibilities, the Bayesian
influence function --- each graded on concept- and
computation-maturity, at the reader's scale.

\item \textbf{The horizon.}
Phase transitions and quasi-singular phenomena as delivered
loss-landscape consequences; the specification gap as a research
program --- if hyperparameters do implicit geometric selection,
principled hyperparameters become a well-posed question, with an
ancestor in the literature; and the open unification question, with
one developed flag (interpretability) and two credible sketches
(activations, policies).

\end{enumerate}

% ============================================================================
\section{What the Paper Is NOT}
% ============================================================================

\begin{itemize}[itemsep=4pt]
\item A claim that ``deep learning is singular'' is novel. The
      contribution is installing the geometry --- entered through the
      set --- as a practitioner-facing object, the lens, the
      two-questions recast, and the state-of-the-field report.
\item An indictment of practice: the repairs succeeded at their
      actual goal (stability); cross-validation is vindicated at its
      job; Adam's guarantee is correct; the tools work through a
      geometry their textbook stories miss.
\item A partition of tools into safe and broken: every tool answers
      some question validly; every tool's valid scope is narrower than
      its routine use. The severity is not uniform: some tools are
      wrong at their own stated job (the alarm); the rest are right at
      a narrower job than they are given.
\item A claim that your optimizer is broken: the claim is that its
      choices are currently unexplained, and newly measurable.
\item A claim that singular tools dominate everywhere: nothing here
      improves on Adam for very large models, and that is fine; it is
      a different axis.
\item A claim that the reader can algorithmically classify any given
      failure: the lens sets expectations and tags tools; it is not a
      diagnostic procedure.
\item A claim of full rigor on ReLU: the theorems require
      analyticity; the paper says so plainly and works where the
      mathematics holds.
\item A textbook or proof-driven treatment of SLT.
\item A survey of loss-landscape, symmetry, or representation-learning
      literature.
\item An argument for Bayesian over frequentist deep learning.
\item A claim that singular geometry closes double descent or scaling:
      it makes them askable.
\item A purely theoretical paper: it recasts real tools and reports
      honest maturity for those that exist.
\end{itemize}

% ============================================================================
\section{Paper Structure Overview}
% ============================================================================

Each section is described by its \textbf{goal}, \textbf{main
discussion}, and \textbf{close}. This is a writing guide, not a table of
contents. The set threads the spine without reorganizing it: assumed
away in \S1, installed in \S2, formalized in \S3, read and chosen from
in \S4, opened as a program in \S5.

\bigskip

\noindent\textbf{\S1. Something Is Wrong} ($\sim$3 pages)

\textbf{Goal:} Urgency, and only urgency --- and only level-1 urgency:
tools that make claims in parameter space and are wrong at their own
stated jobs. The governance story (level 2) is planted as a question,
not argued here.

\textbf{Main discussion:} Four escalating pillars: wrong results (VI,
normalization's rationale) $\to$ misaimed fixes (damping,
sharpness-seeking, isotropizing) $\to$ regularizing instead of
redesigning $\to$ folklore in the gap. The pillar-tool audit table is
the centerpiece, with a third gate added: \emph{every row makes a
parameter-space claim} (posterior, uncertainty region, influence,
curvature, parameter count). The convergence-guarantees row moves out of
the table --- it is a level-2 item and lives in \S4.2's recast. The
masking thesis stated as the common thread, with its honest history: the
repairs were aimed at stability and delivered it; the failure is a
byproduct of a success, which is why no one noticed and no one is
accused. Cross-validation spared, deferred to \S4.

\textbf{Close:} One cause, silently masked; a question the reader cannot
yet ask precisely (what did my optimizer choose, and from what?); why
does it all fail at once? Hand off to \S2. Source for the companion
blog.

\bigskip

\noindent\textbf{\S2. Communicating the Singular Geometry} ($\sim$4
pages)

\textbf{Goal:} Answer \S1's question by installing the object --- the
correct geometric picture without algebraic-geometry language.

\textbf{Main discussion:} The two engines (symmetry via normalization,
$f(x,\alpha\theta)=f(x,\theta)$; overparameterization). Normal crossing,
not wide valley --- hyperbolic level sets as the empirical signature.
Underdetermination stated where the picture makes it obvious: every
point of the crossing is loss-equivalent, so ``minimize the loss'' does
not finish the job; something else chooses, and \S4 names the
instrument that reads the choice. The descent in visibility on the
microscope FFN: plot the crossing, calibrate the probes, climb in
dimension until only the probes survive. Then one argument in two
halves: the high-dimensionality reckoning (the curse of dimensionality
forecloses classical inspection --- profile likelihood, visualized
posteriors, sampler-resolved level sets --- which is the argument for
probes and ultimately SLT) and the posterior pinned to the singular set,
evolving with $n$ (the sample-size arm made visible: the level sets fix
how much room finite-$n$ estimates have to wander; growing $n$ closes
the room, revealing the set --- and the compounding stated: the set is
revealed by $n$ exactly where dimension makes it expensive to confront
[contingent: C4]). Why standard fixes cannot work (reparameterization,
regularization, healthy-looking diagnostics).

\textbf{Close:} Correct working intuition and a fork: tools (\S4) or
mathematics (\S3). The blow-up is \S3's payload, not a cliffhanger.

\bigskip

\noindent\textbf{\S3. Singular Learning Theory: The Core Framework}
($\sim$2.5--3 pages)

\textbf{Goal:} Formalize \S2's picture. Skippable by the tools-only
practitioner; a from-scratch on-ramp for the theorist without algebraic
geometry.

\textbf{Main discussion:} Regular geometry formalized (Fisher, Laplace,
$p/2$). Singular geometry formalized (the optimal set as a variety;
Fisher degeneracy; the stratified parameter space). The analyticity
caveat, stated plainly: the theorems require analytic models; the
microscope FFN uses an analytic activation so every step is carried out
in full rigor; ReLU is the limit of analytic units and the formal
extension is open. The blow-up: resolution to normal-crossing form;
$\lambda$ emerges; the free-energy expansion. The sorting principle:
integral-based quantities survive the non-smooth change of variables,
local-differential ones do not --- the formal ground of the geometric
arm of the lens (a geometry problem is one where the local-differential
object was never the right object).

\textbf{Close:} Regular geometry is the Laplace picture; singular
geometry breaks it; the blow-up restores tractability; complexity is
governed by the geometry of the minimizing set, not the parameter
count.

\bigskip

\noindent\textbf{\S4. Now What} ($\sim$4--5 pages)

\textbf{Goal:} The delivered value. A maturity-aware field report
organized around the questions tools answer, each instrument graded on
concept- and computation-maturity separately --- computability judged at
the reader's scale.

\textbf{Main discussion:} Practitioner layer. Two questions, two
criteria as the organizing recast: if the question is prediction,
CV/WAIC is the right object and is vindicated in singular models
(Watanabe 2010, gen$+$CV $\to 2\lambda/n$; AIC/DIC provably fail); if
the question is identifying structure, the free energy is the right
object, WBIC estimates it, and $\lambda$-based selection is consistent
where CV is not --- singularity lets CV tolerate larger models (smaller
$\lambda$), and at small $n$ the selection signal shrinks while CV's
noise does not [contingent: C3]. Adam recast through underdetermination:
the guarantee is honest (an $\epsilon$-stationary point under
smoothness) and the target is not a well-posed destination --- the
objective specifies a set; the hyperparameters silently complete the
choice; whether $\epsilon$ and its kin merely navigate the degenerate
landscape or select the stratum is posed as the open question it is,
with Adam's $\epsilon$ flagged as a candidate damping instance
[VERIFY]. The tool-by-question scope table (what it answers / what
practice asks / the gap), with safe/caution/avoid as its practitioner
projection. Stop here if you want.

Methods layer. The LLC as the readout of where training landed: the
SGD-vs-entropy-SGD exhibit (1.9M-parameter FFN on MNIST; 80 repetitions;
entropy-SGD lands on measurably lower-$\lambda$ strata with remarkably
low variance; matched fit [VERIFY, C7]) --- optimizer inductive bias
made visible by the number the loss cannot see; NGD-vs-SGD as the second
instance. Singular VI as the constructive fix (the \S1 indictment
becomes the tool). The estimation toolkit (LLC/rLLC, \texttt{devinterp},
susceptibilities). The Bayesian influence function (the damping constant
returns as the well-posed localization $\gamma$ --- the twin arc).
Sampler craft on the singular strata as the live computational frontier.
The worked example (rLLC on attention heads) in the appendix with a
one-paragraph pointer, bridging to \S5's interpretability flag.

\textbf{Close:} Which tools to trust for which question, which are
real-but-early, and which kind of hardness each addresses. The good
instincts were right --- and every one of them, including the spared
tool, is right under a scope. The delivered value closes here.

\bigskip

\noindent\textbf{\S5. Conclusions and the Open Horizon} ($\sim$2--2.5
pages)

\textbf{Goal:} Conclude the delivered paper; pose the open question;
plant a substantive but explicitly conjectural flag.

\textbf{Main discussion:} Recap (the loss-landscape case is solid).
Delivered consequences: phase transitions (regularizers select strata; a
mechanism for part of double descent), quasi-singular phenomena (RLCT
growth; bearing on scaling). The specification gap as a program:
hyperparameters silently complete a choice the objective cannot express;
the LLC makes the choice measurable; an optimizer targeting the free
energy already demonstrates explicit selection --- so principled, even
optimal, hyperparameters are a well-posed open question (explicitly
conjectural; hyperparameter-dependence of final $\lambda$ within a
single algorithm flagged as an open experiment). The open question:
could singular geometry unify deep learning practice? The flag:
mechanistic interpretability developed at half-page depth (non-unique
circuits, superposition as degeneracy on the singular set, the rLLC
bridge from \S4's appendix); activation/representation space and
policy/value space compressed to a credible sentence each. Art to
science: the \S4 sampler frontier points toward geometry-preserving
machinery of the kind physics already has for singular configuration
spaces; the ReLU extension named alongside it as open mathematics.

\textbf{Close:} The loss landscape is governed by singular geometry and
there are tools to use now. The unification is open --- the case is
strong enough to pursue. Invitation, not synthesis.

\bigskip
\noindent Total: $\sim$16--17 pages, plus mathematical appendix (theorem
statements; the worked rLLC example). Practitioner-first budget: \S1,
\S2, \S4 protected; \S3 and \S5 compress first.

% ============================================================================
\section{Key Framing Decisions}
% ============================================================================

\begin{enumerate}[label=(\alph*), itemsep=6pt]

\item \textbf{The set is the door; the geometry is the subject.} The
paper's entry move is a substitution --- the minimum is a set, not a
point --- because it is the one sentence a practitioner absorbs
instantly (``you did not find the minimum; there is no `the.'\,'').
But the set alone determines almost nothing the paper delivers:
$z_1^2$ and $z_1^4$ share a solution set and differ in $\lambda$,
posterior, and free energy. SLT's quantities are invariants of the
\emph{pair} --- the set and how the loss vanishes along it --- and the
vanishing structure is where every consequence lives. The install
sequence, which is the spine: the set (why it is there) $\to$ the
shape of the loss around it (normal crossing) $\to$ the standard form
and the invariants ($\lambda$, singular fluctuation) $\to$ the
posterior and free energy (how $n$ enters) $\to$ the consequences for
training, evaluation, selection, and architecture. Masking, the lens,
scope-narrowing, underdetermination, the instruments --- all
consequences of taking the geometry seriously, entered through the
set. \emph{Note: in this document, ``the set'' abbreviates the pair;
the paper's prose must not inherit the abbreviation where the
vanishing structure is doing the work.}

\item \textbf{The spine is problem $\to$ mechanism $\to$ resolution
$\to$ conclusions $+$ horizon.} Unification is a question \S5 raises,
not a result; delivered value closes at \S4.

\item \textbf{The alarm is level 1; the governance story is level 2.}
Level 1: tools making parameter-space claims falsely assume the set is a
point and are wrong at their own stated jobs --- this, and only this, is
\S1's alarm, and it gives the audit table its third gate (every row
makes a parameter-space claim). Level 2: procedures whose stated jobs
survive but whose outcomes the set silently governs --- which point the
optimizer selects, why it generalizes, which model is chosen at finite
$n$. Level 2 carries no indictment; its register is explanatory and its
payoff is \S4's instruments and \S5's program. The bridge, stated once
and plainly: the optimizer's silent choice has generalization
consequences ($\lambda$ enters the generalization error), so the
geometry reaches what the practitioner cares about even when every
number on the screen looks healthy.

\item \textbf{The harm statement is masking, honestly historied.} The
repairs were aimed at stability, and stability is what they delivered
--- the masking is a byproduct of a success, not a failed attempt at
removing the geometry. Every regularity-restoring repair returns a
plausible number, so the failure is silent: structurally wrong where
checkable, distorted elsewhere. This is why the alarm accuses no one.
``Masked, not removed'' is retired from the opening and survives
downstream (\S1's close, \S2.5) where it describes what the repairs do,
not what anyone intended.

\item \textbf{Underdetermination is the optimizer-facing face of the
object; scope-narrowing is the evaluative-facing face.} Siblings, both
children of the set. Minimization on a singular landscape is an
incomplete specification: the loss cannot prefer among loss-equivalent
points, so the optimizer and hyperparameters complete the choice ---
Adam is not wrong at its job; its job is underdetermined, and the tuning
folklore is residual selection done by hand, without a name (Pillar 4's
deepest face: folklore accumulates in a \emph{specification} gap, not
just a knowledge gap). For evaluative tools the move is scope-narrowing:
every number answers a narrower question than practice asks of it. Twin
calibrations: Adam (the guarantee certifies a stationary point, nothing
about which) and CV (a generalization estimator, inconsistent for
selection). When the workhorse and the spared tool submit to the same
move, the move is the thesis's operational form.

\item \textbf{The lens is three-way and threads the spine.} Geometric
hardness (wrong object) / sample size (unrevealed set) / computation
(right object, expensive in high dimension). Any may be operative; they
co-occur and compound: growing $n$ reveals the set exactly where
dimension makes it expensive to confront --- the DL-specific enrichment
the applied-stats sibling never needed [contingent: C4]. Surfaced as
confusion in \S1, made recognizable in \S2--3, cashed out per tool in
\S4, placed as frontier in \S5. Expectation-setting and per-tool
tagging, not a diagnostic procedure; the operational corollary is the
scope table, and the daily habit is the scope question: which question
does this number actually answer?

\item \textbf{Seen but not mainstreamed.} The information-geometry
school (Fukumizu, Amari) identified the structure; Watanabe supplied the
algebraic-geometric machinery; the observational literature (mode
connectivity, Hessian spectra, re-basin) mapped the set without
receiving it as an object. Context, not the goal: it makes the
counter-evidence the paper's own evidence and keeps the tone generous.
The lineage stays institutional, not biographical.

\item \textbf{Two questions, two criteria (the hinge, rescoped).} CV is
vindicated \emph{as a generalization estimator}: asymptotically
equivalent to WAIC in singular models, gen$+$CV $\to 2\lambda/n$, where
AIC/DIC provably fail. For the selection question, the free energy is
the right object and is consistent where CV/WAIC is not; Watanabe states
both halves and himself deferred the selection comparison. ``CV
vindicated'' must never slide into ``CV suffices'' --- even the spared
tool is spared under a scope, which makes the paper's move uniform with
no exceptions. Bonus the lens predicts: the first place a singular tool
beats the spared classical one at its own job ($\lambda$-based
selection), register subject to decision (j).

\item \textbf{The register is assigned by the audit, not by
enthusiasm.} Plainly provocative only where the tag is geometric
\emph{and} a better tool is computable at the reader's scale today;
equally plain and humble where the geometry is settled but computation
is frontier. Scale is a column in the \S4 audit, and the audit decides
which register each tool earns. Uniform sharpness overclaims; uniform
humility loses the edge.

\item \textbf{Three reader types, three paths.} Practitioner: \S1 $\to$
\S2 $\to$ \S4. Theorist without algebraic geometry: \S1 $\to$ \S2 $\to$
\S3 $\to$ \S4. Contributor: through \S5. Practitioner-first in center of
gravity, and the primary reader is defined by product: a supported
claim, not just a running system.

\item \textbf{Three model roles.} The microscope (tiny FFN with an
analytic activation, \S1--2: make the crossing visible and the
mathematics exact); the reckoning (\S2: visualization fails in real
models --- itself the argument for probes and SLT); the workbench (real
transformer, real BNN, \S4: demonstrate tools where computation-maturity
problems are real).

\item \textbf{Four pillars, evidence calibrated per pillar, gated to
level 1.} Demonstrated $\to$ demonstrable $\to$ argued from the pattern
$\to$ argued from vetted instances. The ordering is the argument; the
audit table is the centerpiece; every row makes a parameter-space claim.
The convergence-lore folklore face moves from \S1 to \S4 (it is level
2); \S1's folklore faces are the variational-gap misattribution and the
damping ritual.

\item \textbf{Normal crossing, not wide valley.} \S2's central
mental-model correction; sets up the blow-up as the obvious untangling.

\item \textbf{Integral-based survives; local-differential fails.} The
sorting principle that makes the alarm precise, spares cross-validation,
and grounds the geometric arm of the lens. Stated in \S3, spent in \S4.

\item \textbf{\S4 is a field report in two layers, written from
audits, not impressions.} Practitioner (which question, which tool) and
methods (how to compute, failure modes); concept-maturity and
computation-maturity graded separately; \texttt{devinterp} the methods
anchor, real-but-early. Two gating tasks: the tools-and-maturity audit
(per tool: problem / tag / maturities / scale / rough edges / layer) and
the practice inventory (per practitioner activity: claim made / tool
used / where the computation lives / the question it answers vs the
question asked / verdict). \S4 gets the pages it needs; \S3 and \S5
compress first.

\item \textbf{The singular-first practice is proof of concept.}
\emph{Deep Learning is Singular, and That's Good} (the invitation);
Lau et al.'s LLC (the instrument, with the SGD-vs-entropy-SGD exhibit);
Wei--Lau singular VI; susceptibilities; the Bayesian influence function
--- cited generously as the toolkit being born. Two emblematic exhibits:
the BIF (a folklore damping constant reappearing as a principled,
computable quantity) and the LLC optimizer experiment (equivalent by
every number you currently look at; different by the number you don't
--- pending C7).

\item \textbf{Approachable re-exposition is a contribution.} Watanabe
built the field from the algebra up; this paper starts from phenomena
the reader has met and pulls down only the algebra each needs. Success:
the reader engages with SLT and prefers a singular tool where the
regular one papers over the geometry.

\item \textbf{Minimal prerequisites, maximum geometric intuition.}
Assumed: SGD/Adam, Hessians, weight decay, normalization, transformers,
double descent and mode connectivity as phenomena. Not assumed:
algebraic geometry, differential geometry, SLT, Watanabe.

\item \textbf{Tone: direct, generous, calibrated.} Not ``you have been
doing it wrong''; not ``singular geometry solves everything.'' The
opening keeps one clause of generosity (the field has seen the set); the
paper includes at least one exhibit of the field working an axis
correctly (statistics acknowledging data batching, \S4) --- generosity
shown, not claimed. ``An object you were never handed'' must never
read as ``an object you were too naive to notice.'' The closing note is
an invitation.

\item \textbf{The analyticity gap is stated, not stepped around.}
Watanabe's theorems require analytic models; ReLU is not analytic. The
microscope FFN uses an analytic activation so every step --- crossing,
blow-up, $\lambda$, validation --- is carried out in full rigor. ReLU is
presented as the limit of analytic units (softplus, $\beta\to\infty$)
whose exact positive homogeneity strengthens the symmetry argument; the
gap is in the machinery, not the phenomenon. Modern transformer
activations (GELU, SiLU) are analytic, narrowing the gap for the
reader's own models [contingent: C6]. Stated as a \S3 caveat; the formal
extension to ReLU flagged as open in \S5. The paper's ethic --- name the
assumption --- applies to the paper.

\item \textbf{Aims are four: alarm, picture $+$ lens, equip $+$
report, and the conjectural extension.} The lens and underdetermination
live inside aim (ii) as the set's two working consequences rather than
as separate procedural promises. The first three aims parallel the
applied-stats sibling exactly; aim (iv) is the DL-specific addition,
gated to the conjectural register by \S5's scope.

\end{enumerate}

% ============================================================================
\section{Load-Bearing Contingencies}
% ============================================================================

\noindent\emph{Framing in this document that depends on an open or
partial item in the needs document. Each entry states the claim, the
needs item it rests on, the decided fallback if verification fails, and
--- where verification purchases something --- the upside action. No
other framing decision is contingent; if one of these fails, only its
listed downstream edits are infected.}

\begin{enumerate}[label=\textbf{C\arabic*.}, leftmargin=*, itemsep=6pt]

\item \textbf{The conclusion-flip exhibit.} The masking thesis asserts
wrong \emph{conclusions}, not just wrong numbers; silence is also what
correctness predicts. Needed: one validated case where a regular tool
and its singular counterpart disagree on a conclusion and the singular
answer is verifiably right. Three candidates, cheapest first: (a) the
small-$n$ mixture selection experiment (CV selects the larger model;
free energy selects the true one) --- doubles as the cross-paper bridge
to the applied-stats binomial mixture; (b) the LLC optimizer exhibit if
matched fit holds (C7); (c) BIF-vs-damped-IF reordering of influential
examples, validated by retraining. \emph{Fallback:} the alarm rests on
the structural-wrongness exhibits (Wei--Lau VI, degenerate-Hessian
Laplace); masking is stated as demonstrated distortion plus
unattributable success; the conclusion-flip is named as the open
empirical question (moves to \S5). \emph{Touches:} aim (i)'s
``structurally wrong where checkable'' (softens to ``can be
structurally wrong''); \S2 walk-away item 3 survives unchanged (masking
restated as demonstrated distortion); \S4-overview \S1's close. [Needs
doc \S4.]

\item \textbf{Wei--Lau precise statements and scale.} Pillar 1's anchor
and \S4.3's provocative register assume the failure/fix statements pin
down as remembered and that singular VI is computable at a scale
relevant to the reader. \emph{Fallback:} if the scale check fails, the
tag becomes geometry-settled / computation-frontier and the register
drops to humble per decision (i); Pillar 1 leans on the Laplace exhibit
as anchor. \emph{Touches:} \S4-overview \S1 and \S4; framing decision
(p). [Needs doc \S1, \S4.]

\item \textbf{The small-$n$ selection statement.} The two-questions
hinge is verified in its qualitative form (CV/WAIC unbiased for
generalization, not consistent; free energy consistent, not unbiased;
the WBIC discussion's weaker-consistency remark; Watanabe's own deferral
of the selection comparison). What remains open is the quantitative
small-$n$ claim (selection signal shrinking against CV noise): pin it in
Watanabe's 2021 mixture comparison or Yamazaki--Kaji 2013, or
demonstrate it with the mixture experiment (C1a). \emph{Fallback:} \S4.1
states the consistency asymmetry without the finite-$n$ rate; the
mixture experiment is reported as illustration, not theorem.
\emph{Touches:} \S4-overview \S4's practitioner layer only; walk-away
item 6 survives unchanged (it states the consistency asymmetry, not the
rate). [Needs doc \S4.]

\item \textbf{The $n$--dimension compounding claim.} The compounding
sentence assumes the claim survives the finite-sample /
dimension-dependence literature check --- which also asks whether the
compounding is genuinely unspelled-out (a contribution to name) or
already stated (cite and lean). \emph{Fallback:} shown-not-interpreted
in \S2 (the figure stands on its own), strong form posed as open in
\S5. \emph{Touches:} \S4-overview \S2.4; framing decision (f)'s
compounding sentence. \emph{If verified:} restore the compounding to
NOT bullet 1's contribution list and to walk-away item 5's lens
sentence. [Needs doc \S2.]

\item \textbf{The one-reflex reading of Pillar 3.} The claim that
damping, Gaussian families, and sharpness penalties are one move at
different sites must hold across the audit table's shipped-repair column
--- now easier, since the table is gated to parameter-space claims.
\emph{Fallback:} Pillar 3 weakens from ``the same move'' to ``a family
of regularity-restoring moves''; the escalation ordering of \S1 survives
intact. \emph{Touches:} \S4-overview \S1 only. [Needs doc \S1.]

\item \textbf{Transformer activations are analytic.} Decision (t)'s
``narrowing the gap'' sentence assumes GELU/SiLU dominance in the
workbench class and analyticity of those units. \emph{Fallback:} the
sentence is cut; the analyticity caveat stands on the microscope's rigor
and the softplus-limit framing alone. \emph{Touches:} framing decision
(t) only. [Needs doc \S3.]

\item \textbf{Matched fit in the LLC optimizer exhibit.} The emblematic
reading (``equivalent by every number you currently look at; different
by the number you don't'') requires that the SGD and entropy-SGD
endpoints have comparable training and test performance in Lau et al.'s
experiment --- the $\lambda$ separation and its low variance are
verified; the matched-fit half is not yet. \emph{Fallback:} the exhibit
is reported as optimizer inductive bias made measurable (verified as
stated by the authors), without the equivalence framing; the emblematic
slot falls to the BIF alone, and C1(b) is struck as a conclusion-flip
candidate. \emph{Touches:} \S4-overview \S4's methods layer; framing
decision (p); walk-away item 6 survives unchanged (it states the LLC
as readout, not the equivalence framing). [Needs doc \S4, new item.]

\end{enumerate}

\end{document}
